\section{Methodology}

\subsection{Research Materials}
The core dataset consists of 416 semantic mutations, constructed by taking real, verified context documents and systematically injecting semantic faults into them. Unlike simple rule-based mutations used in previous work --- such as randomly deleting characters or swapping synonyms --- the semantic mutations in this study directly target informational content: altering factual claims, swapping critical numerical values, or inserting buggy code snippets, all while keeping the text grammatically fluent and contextually coherent.

For each of the 416 mutations, the dataset records four components:

\begin{itemize}
\item The user query (Query).
\item The original context (containing verified information).
\item The baseline ground-truth answer.
\item The mutated context (poisoned).
\end{itemize}

\subsection{Experimental Procedure}
For each experimental configuration, the corresponding text generation model (Llama 3.1 8B or 17B) was given the user query along with the mutated context as input. The model-generated responses (Mutated Responses) were then collected and compiled into CSV datasets for the evaluation step.

\subsection{Evaluation Pipeline}
Evaluating 416 responses across multiple configurations requires a process that is both rigorous and scalable. We therefore adopted a hybrid AI-human evaluation approach with two main components:

\begin{itemize}
\item \textbf{AI Judge (Primary Evaluator)}: The model \textit{meta-llama/llama-4-scout} (17B) was tasked with evaluating all 416 RAG responses for each configuration. The AI Judge was provided with the query, the mutated context, the ground-truth answer, and the RAG-generated response when given the mutated context. It then classified each response into one of four mutually exclusive labels:
\begin{enumerate}[label=(\alph*)]
\item \textit{Faithful}: The model passively follows the mutated context, repeating the misinformation without any intervention from external knowledge.
\item \textit{Inconsistent}: The model recognizes the discrepancy in the context and uses its internal parametric knowledge to self-correct, producing a factually accurate answer that does not match the mutated text.
\item \textit{Abstain}: The model explicitly refuses to answer, citing reasons related to contradictory or insufficient information.
\item \textit{Hallucination}: The model fabricates information that appears in neither the mutated context nor the ground-truth answer.
\end{enumerate}
\item \textbf{Human Evaluator}: To verify the accuracy of the AI Judge's classifications, we also had a human evaluator independently score a fixed, randomly selected subset of 42 samples per configuration. The human evaluations serve as the ground-truth benchmark for measuring the reliability of the AI Judge.
\end{itemize}

\subsection{Data Analysis}
Statistical rigor is a key concern in our methodology. To validate the AI Judge's reliability, we computed Cohen's Kappa ($k$) to measure the level of agreement between the AI Judge and the human evaluator on the 42-sample subsets. We chose Cohen's Kappa because it accounts for the possibility that agreement could happen by chance, making it a more trustworthy measure of inter-rater reliability compared to simple agreement percentages.

Additionally, to test whether the RAG system can effectively maintain a defensive stance, we set up a null hypothesis ($H_0$) stating that a safe RAG system should achieve an Abstain rate of at least $\ge 90\%$ when faced with semantic mutations. We then applied exact Binomial tests (using \texttt{scipy.stats.binomtest}) to determine whether the observed Abstain rates were statistically significant relative to this threshold.
